% Options for packages loaded elsewhere
% Options for packages loaded elsewhere
\PassOptionsToPackage{unicode}{hyperref}
\PassOptionsToPackage{hyphens}{url}
\PassOptionsToPackage{dvipsnames,svgnames,x11names}{xcolor}
%
\documentclass[
]{article}
\usepackage{xcolor}
\usepackage[top=20mm,left=20mm,right=20mm,heightrounded]{geometry}
\usepackage{amsmath,amssymb}
\setcounter{secnumdepth}{-\maxdimen} % remove section numbering
\usepackage{iftex}
\ifPDFTeX
  \usepackage[T1]{fontenc}
  \usepackage[utf8]{inputenc}
  \usepackage{textcomp} % provide euro and other symbols
\else % if luatex or xetex
  \usepackage{unicode-math} % this also loads fontspec
  \defaultfontfeatures{Scale=MatchLowercase}
  \defaultfontfeatures[\rmfamily]{Ligatures=TeX,Scale=1}
\fi
\usepackage{lmodern}
\ifPDFTeX\else
  % xetex/luatex font selection
\fi
% Use upquote if available, for straight quotes in verbatim environments
\IfFileExists{upquote.sty}{\usepackage{upquote}}{}
\IfFileExists{microtype.sty}{% use microtype if available
  \usepackage[]{microtype}
  \UseMicrotypeSet[protrusion]{basicmath} % disable protrusion for tt fonts
}{}
\makeatletter
\@ifundefined{KOMAClassName}{% if non-KOMA class
  \IfFileExists{parskip.sty}{%
    \usepackage{parskip}
  }{% else
    \setlength{\parindent}{0pt}
    \setlength{\parskip}{6pt plus 2pt minus 1pt}}
}{% if KOMA class
  \KOMAoptions{parskip=half}}
\makeatother
% Make \paragraph and \subparagraph free-standing
\makeatletter
\ifx\paragraph\undefined\else
  \let\oldparagraph\paragraph
  \renewcommand{\paragraph}{
    \@ifstar
      \xxxParagraphStar
      \xxxParagraphNoStar
  }
  \newcommand{\xxxParagraphStar}[1]{\oldparagraph*{#1}\mbox{}}
  \newcommand{\xxxParagraphNoStar}[1]{\oldparagraph{#1}\mbox{}}
\fi
\ifx\subparagraph\undefined\else
  \let\oldsubparagraph\subparagraph
  \renewcommand{\subparagraph}{
    \@ifstar
      \xxxSubParagraphStar
      \xxxSubParagraphNoStar
  }
  \newcommand{\xxxSubParagraphStar}[1]{\oldsubparagraph*{#1}\mbox{}}
  \newcommand{\xxxSubParagraphNoStar}[1]{\oldsubparagraph{#1}\mbox{}}
\fi
\makeatother

\usepackage{color}
\usepackage{fancyvrb}
\newcommand{\VerbBar}{|}
\newcommand{\VERB}{\Verb[commandchars=\\\{\}]}
\DefineVerbatimEnvironment{Highlighting}{Verbatim}{commandchars=\\\{\}}
% Add ',fontsize=\small' for more characters per line
\usepackage{framed}
\definecolor{shadecolor}{RGB}{241,243,245}
\newenvironment{Shaded}{\begin{snugshade}}{\end{snugshade}}
\newcommand{\AlertTok}[1]{\textcolor[rgb]{0.68,0.00,0.00}{#1}}
\newcommand{\AnnotationTok}[1]{\textcolor[rgb]{0.37,0.37,0.37}{#1}}
\newcommand{\AttributeTok}[1]{\textcolor[rgb]{0.40,0.45,0.13}{#1}}
\newcommand{\BaseNTok}[1]{\textcolor[rgb]{0.68,0.00,0.00}{#1}}
\newcommand{\BuiltInTok}[1]{\textcolor[rgb]{0.00,0.23,0.31}{#1}}
\newcommand{\CharTok}[1]{\textcolor[rgb]{0.13,0.47,0.30}{#1}}
\newcommand{\CommentTok}[1]{\textcolor[rgb]{0.37,0.37,0.37}{#1}}
\newcommand{\CommentVarTok}[1]{\textcolor[rgb]{0.37,0.37,0.37}{\textit{#1}}}
\newcommand{\ConstantTok}[1]{\textcolor[rgb]{0.56,0.35,0.01}{#1}}
\newcommand{\ControlFlowTok}[1]{\textcolor[rgb]{0.00,0.23,0.31}{\textbf{#1}}}
\newcommand{\DataTypeTok}[1]{\textcolor[rgb]{0.68,0.00,0.00}{#1}}
\newcommand{\DecValTok}[1]{\textcolor[rgb]{0.68,0.00,0.00}{#1}}
\newcommand{\DocumentationTok}[1]{\textcolor[rgb]{0.37,0.37,0.37}{\textit{#1}}}
\newcommand{\ErrorTok}[1]{\textcolor[rgb]{0.68,0.00,0.00}{#1}}
\newcommand{\ExtensionTok}[1]{\textcolor[rgb]{0.00,0.23,0.31}{#1}}
\newcommand{\FloatTok}[1]{\textcolor[rgb]{0.68,0.00,0.00}{#1}}
\newcommand{\FunctionTok}[1]{\textcolor[rgb]{0.28,0.35,0.67}{#1}}
\newcommand{\ImportTok}[1]{\textcolor[rgb]{0.00,0.46,0.62}{#1}}
\newcommand{\InformationTok}[1]{\textcolor[rgb]{0.37,0.37,0.37}{#1}}
\newcommand{\KeywordTok}[1]{\textcolor[rgb]{0.00,0.23,0.31}{\textbf{#1}}}
\newcommand{\NormalTok}[1]{\textcolor[rgb]{0.00,0.23,0.31}{#1}}
\newcommand{\OperatorTok}[1]{\textcolor[rgb]{0.37,0.37,0.37}{#1}}
\newcommand{\OtherTok}[1]{\textcolor[rgb]{0.00,0.23,0.31}{#1}}
\newcommand{\PreprocessorTok}[1]{\textcolor[rgb]{0.68,0.00,0.00}{#1}}
\newcommand{\RegionMarkerTok}[1]{\textcolor[rgb]{0.00,0.23,0.31}{#1}}
\newcommand{\SpecialCharTok}[1]{\textcolor[rgb]{0.37,0.37,0.37}{#1}}
\newcommand{\SpecialStringTok}[1]{\textcolor[rgb]{0.13,0.47,0.30}{#1}}
\newcommand{\StringTok}[1]{\textcolor[rgb]{0.13,0.47,0.30}{#1}}
\newcommand{\VariableTok}[1]{\textcolor[rgb]{0.07,0.07,0.07}{#1}}
\newcommand{\VerbatimStringTok}[1]{\textcolor[rgb]{0.13,0.47,0.30}{#1}}
\newcommand{\WarningTok}[1]{\textcolor[rgb]{0.37,0.37,0.37}{\textit{#1}}}

\usepackage{longtable,booktabs,array}
\newcounter{none} % for unnumbered tables
\usepackage{calc} % for calculating minipage widths
% Correct order of tables after \paragraph or \subparagraph
\usepackage{etoolbox}
\makeatletter
\patchcmd\longtable{\par}{\if@noskipsec\mbox{}\fi\par}{}{}
\makeatother
% Allow footnotes in longtable head/foot
\IfFileExists{footnotehyper.sty}{\usepackage{footnotehyper}}{\usepackage{footnote}}
\makesavenoteenv{longtable}
\usepackage{graphicx}
\makeatletter
\newsavebox\pandoc@box
\newcommand*\pandocbounded[1]{% scales image to fit in text height/width
  \sbox\pandoc@box{#1}%
  \Gscale@div\@tempa{\textheight}{\dimexpr\ht\pandoc@box+\dp\pandoc@box\relax}%
  \Gscale@div\@tempb{\linewidth}{\wd\pandoc@box}%
  \ifdim\@tempb\p@<\@tempa\p@\let\@tempa\@tempb\fi% select the smaller of both
  \ifdim\@tempa\p@<\p@\scalebox{\@tempa}{\usebox\pandoc@box}%
  \else\usebox{\pandoc@box}%
  \fi%
}
% Set default figure placement to htbp
\def\fps@figure{htbp}
\makeatother





\setlength{\emergencystretch}{3em} % prevent overfull lines

\providecommand{\tightlist}{%
  \setlength{\itemsep}{0pt}\setlength{\parskip}{0pt}}



 


\makeatletter
\@ifpackageloaded{tcolorbox}{}{\usepackage[skins,breakable]{tcolorbox}}
\@ifpackageloaded{fontawesome5}{}{\usepackage{fontawesome5}}
\definecolor{quarto-callout-color}{HTML}{909090}
\definecolor{quarto-callout-note-color}{HTML}{0758E5}
\definecolor{quarto-callout-important-color}{HTML}{CC1914}
\definecolor{quarto-callout-warning-color}{HTML}{EB9113}
\definecolor{quarto-callout-tip-color}{HTML}{00A047}
\definecolor{quarto-callout-caution-color}{HTML}{FC5300}
\definecolor{quarto-callout-color-frame}{HTML}{acacac}
\definecolor{quarto-callout-note-color-frame}{HTML}{4582ec}
\definecolor{quarto-callout-important-color-frame}{HTML}{d9534f}
\definecolor{quarto-callout-warning-color-frame}{HTML}{f0ad4e}
\definecolor{quarto-callout-tip-color-frame}{HTML}{02b875}
\definecolor{quarto-callout-caution-color-frame}{HTML}{fd7e14}
\makeatother
\makeatletter
\@ifpackageloaded{caption}{}{\usepackage{caption}}
\AtBeginDocument{%
\ifdefined\contentsname
  \renewcommand*\contentsname{Table of contents}
\else
  \newcommand\contentsname{Table of contents}
\fi
\ifdefined\listfigurename
  \renewcommand*\listfigurename{List of Figures}
\else
  \newcommand\listfigurename{List of Figures}
\fi
\ifdefined\listtablename
  \renewcommand*\listtablename{List of Tables}
\else
  \newcommand\listtablename{List of Tables}
\fi
\ifdefined\figurename
  \renewcommand*\figurename{Figure}
\else
  \newcommand\figurename{Figure}
\fi
\ifdefined\tablename
  \renewcommand*\tablename{Table}
\else
  \newcommand\tablename{Table}
\fi
}
\@ifpackageloaded{float}{}{\usepackage{float}}
\floatstyle{ruled}
\@ifundefined{c@chapter}{\newfloat{codelisting}{h}{lop}}{\newfloat{codelisting}{h}{lop}[chapter]}
\floatname{codelisting}{Listing}
\newcommand*\listoflistings{\listof{codelisting}{List of Listings}}
\makeatother
\makeatletter
\makeatother
\makeatletter
\@ifpackageloaded{caption}{}{\usepackage{caption}}
\@ifpackageloaded{subcaption}{}{\usepackage{subcaption}}
\makeatother
\usepackage{bookmark}
\IfFileExists{xurl.sty}{\usepackage{xurl}}{} % add URL line breaks if available
\urlstyle{same}
\hypersetup{
  pdftitle={LA-2: Setting up R (10 points)},
  colorlinks=true,
  linkcolor={blue},
  filecolor={Maroon},
  citecolor={Blue},
  urlcolor={blue},
  pdfcreator={LaTeX via pandoc}}


\title{LA-2: Setting up R (10 points)}
\author{}
\date{}
\begin{document}
\maketitle


\section{Learning Outcomes}\label{learning-outcomes}

In this assignment, you will learn how to:

\begin{itemize}
\tightlist
\item
  Setup an account on Posit Cloud.
\item
  Setup Your Workspace on Posit Cloud.
\item
  Setup a Project on Posit Cloud.
\item
  Write a function in R.
\item
  Create objects in R.
\item
  Call objects that you have created.
\item
  Use R as a calculator.
\end{itemize}

\begin{center}\rule{0.5\linewidth}{0.5pt}\end{center}

We will be using Posit Cloud for lab assignments in this course.
RStudio, which is the graphical user interface through which many people
use R, re-branded themselves as Posit. What used to be called RStudio
Cloud is now known as Posit Cloud, but it is the same thing.

You can also use R and RStudio as standalone programs offline on your
computer, but you will have to download both programs and install them.
We will not be covering installation of R on your own machine in this
course, but you are welcome to do so on your own.

\begin{tcolorbox}[enhanced jigsaw, arc=.35mm, bottomrule=.15mm, bottomtitle=1mm, breakable, colback=white, colbacktitle=quarto-callout-tip-color!10!white, colframe=quarto-callout-tip-color-frame, coltitle=black, left=2mm, leftrule=.75mm, opacityback=0, opacitybacktitle=0.6, rightrule=.15mm, title=\textcolor{quarto-callout-tip-color}{\faLightbulb}\hspace{0.5em}{Tip}, titlerule=0mm, toprule=.15mm, toptitle=1mm]

Remember to save your work often.

\end{tcolorbox}

\section{Setting up an account on Posit
Cloud}\label{setting-up-an-account-on-posit-cloud}

\begin{enumerate}
\def\labelenumi{\arabic{enumi})}
\item
  Open a web browser and navigate to
  \href{https://posit.cloud/}{Posit.Cloud}.
\item
  Click on \texttt{Log\ In} in the top right corner of the page.
\item
  On the next screen, click on the \texttt{Sign\ Up} tab or sign in with
  your Google or Github account (pick one if you have it and would like
  to use it for the course). You may sign up with any account that you
  would like.
\item
  Once you have created an account, log in.
\end{enumerate}

\section{Setting up Posit Cloud}\label{setting-up-posit-cloud}

\begin{enumerate}
\def\labelenumi{\arabic{enumi})}
\item
  Once you log in, you will be taken to \texttt{Your\ Workspace}.
\item
  Click on \texttt{New\ Project} in the top right corner. From the drop
  down options, select \texttt{New\ RStudio\ Project}.
\item
  At the top of the screen, click on \texttt{Untitled\ Project} to name
  your project. I recommend having a project for each lab assignment.
\item
  You're ready to work with R on Posit Cloud! Your instructor will help
  you get familiar with the layout of R. Be sure you know where the
  \textbf{Console} and \textbf{Environment} panes are located.
\item
  Take a screenshot of your project and save it. This is one of your
  submissions for this assignment.
\end{enumerate}

\section{Your First Functions}\label{your-first-functions}

At its core, R is a calculator. It can also make fancy figures. Let's
start by giving R some simple commands.

\begin{enumerate}
\def\labelenumi{\arabic{enumi})}
\item
  Create a new R script, name it \texttt{LA-2\_Your\ Name.R}. Your
  instructor will show you how. Think of a script as a list of commands
  that you (the user) are giving to the computer (R, which is the engine
  underlying RStudio and Posit Cloud).
\item
  In your script, type the following command:
\end{enumerate}

\begin{Shaded}
\begin{Highlighting}[]
\FunctionTok{print}\NormalTok{(}\StringTok{"Your Name Here"}\NormalTok{)}
\end{Highlighting}
\end{Shaded}

\begin{enumerate}
\def\labelenumi{\arabic{enumi})}
\setcounter{enumi}{2}
\item
  Run the command you just wrote in the script and make note of what
  happens in the Console. To run commands from the script, highlight the
  command you want to run and click \texttt{Run} in the top right of the
  script or use the shortcut keys, \texttt{Ctrl} + \texttt{Enter}
  (Windows) or \texttt{Command} + \texttt{Enter}.
\item
  Copy and paste the output from the Console into your R script then add
  \texttt{\#} in front of it.
\end{enumerate}

\begin{tcolorbox}[enhanced jigsaw, arc=.35mm, bottomrule=.15mm, bottomtitle=1mm, breakable, colback=white, colbacktitle=quarto-callout-note-color!10!white, colframe=quarto-callout-note-color-frame, coltitle=black, left=2mm, leftrule=.75mm, opacityback=0, opacitybacktitle=0.6, rightrule=.15mm, title=\textcolor{quarto-callout-note-color}{\faInfo}\hspace{0.5em}{Note}, titlerule=0mm, toprule=.15mm, toptitle=1mm]

Congratulations, you have just programmed your first function in R!

\end{tcolorbox}

With your answer, the script should look something like this:

\begin{Shaded}
\begin{Highlighting}[]
\FunctionTok{print}\NormalTok{(}\StringTok{"Your Name Here"}\NormalTok{)}
\CommentTok{\# [1] "Your Name Here"}
\end{Highlighting}
\end{Shaded}

\begin{enumerate}
\def\labelenumi{\arabic{enumi})}
\setcounter{enumi}{4}
\tightlist
\item
  You can also do basic math with R. In your R script, type the
  following:
\end{enumerate}

\begin{Shaded}
\begin{Highlighting}[]
\DecValTok{5}\SpecialCharTok{*}\DecValTok{12}
\end{Highlighting}
\end{Shaded}

\begin{enumerate}
\def\labelenumi{\arabic{enumi})}
\setcounter{enumi}{5}
\tightlist
\item
  Copy and paste the result from the Console in your R script. Again,
  add \texttt{\#} in front of it.
\end{enumerate}

\section{Creating Objects}\label{creating-objects}

You might notice that the result of the last calculation you did was not
saved anywhere. What if we want to use the result in another calculation
or operation? To do this, we can save functions and computations into
\emph{objects}. Let's create some objects in R so we can use them in
subsequent calculations.

\begin{enumerate}
\def\labelenumi{\arabic{enumi})}
\tightlist
\item
  Type the following in your script and run the function. Describe what
  happened when you ran the command. You can type your answer under the
  function in your R script. Type \texttt{\#} before your answer so that
  the text turns green, indicating that these are comments in your
  script (vs.~functions). Comments are not evaluated when you run the
  script; only functions are evaluated.
\end{enumerate}

\begin{Shaded}
\begin{Highlighting}[]
\NormalTok{xray }\OtherTok{\textless{}{-}} \DecValTok{4}\SpecialCharTok{+}\DecValTok{5}
\end{Highlighting}
\end{Shaded}

\begin{enumerate}
\def\labelenumi{\arabic{enumi})}
\setcounter{enumi}{1}
\item
  In your R script, describe (i) what happened in the Console and (ii)
  what happened in the Environment pane when you ran the function above.
  Remember to use the \texttt{\#} to make all your answers into comments
  in your R script.
\item
  Now, create another object (you can call it whatever you wish) and
  assign the number \texttt{6} to it. Describe what happened in the
  Console and the Environment when you created this object.
\item
  Use the object you just created as the denominator and divide
  \texttt{xray} by this object. Describe what happened in the Console
  and Environment panes when you did this.
\item
  Finally, let's create an object called \texttt{fruitsalad} from
  individual fruit. Create four objects, \texttt{apples},
  \texttt{oranges}, \texttt{pears}, and \texttt{grapes}. Assign the
  following numbers to them: 5, 3, 10, 18, respectively. Create
  \texttt{fruitsalad} by adding the individual fruit together. Write
  down the result of your fruit salad in your R script.
\end{enumerate}

\begin{center}\rule{0.5\linewidth}{0.5pt}\end{center}

\section{Submission}\label{submission}

Submit your R script (named \texttt{LA-\#\_FirstName-LastName.R}) to
Canvas.

Your R script should:

\begin{enumerate}
\def\labelenumi{\arabic{enumi})}
\tightlist
\item
  Include commands and functions that are necessary to address all the
  questions in the assignment.
\item
  Contain comments that answer the questions in the assignment.
\item
  Run in its entirety without errors.
\end{enumerate}

To ensure that your R script runs without errors, you should:

\begin{itemize}
\tightlist
\item
  Save your script.
\item
  Navigate back to Your Workspace on Posit Cloud.
\item
  Reopen your project.
\item
  Run the entire script line-by-line without editing it to ensure there
  are no errors.
\end{itemize}

\begin{tcolorbox}[enhanced jigsaw, arc=.35mm, bottomrule=.15mm, bottomtitle=1mm, breakable, colback=white, colbacktitle=quarto-callout-important-color!10!white, colframe=quarto-callout-important-color-frame, coltitle=black, left=2mm, leftrule=.75mm, opacityback=0, opacitybacktitle=0.6, rightrule=.15mm, title=\textcolor{quarto-callout-important-color}{\faExclamation}\hspace{0.5em}{Important}, titlerule=0mm, toprule=.15mm, toptitle=1mm]

These standards apply to all submissions in this course that require R
scripts. You should follow these instructions for preparation, naming,
and saving of your R script for all of your individual lab assignments.

\end{tcolorbox}




\end{document}
